\documentclass[11pt,letterpaper]{article}

% --- Paquetes de configuración básica ---
\usepackage[utf8]{inputenc}
\usepackage[T1]{fontenc}
\usepackage[spanish,es-noshorthands,es-tabla]{babel}
\usepackage[margin=2cm,top=2.2cm,bottom=2.2cm]{geometry}
\usepackage{amsmath,amssymb}
\usepackage{enumitem}
\usepackage{fancyhdr}
\usepackage{listings}
\usepackage{xcolor}
\usepackage{tcolorbox}
\usepackage{booktabs}
\usepackage{array}

\raggedbottom

% --- Configuración de cabeceras y pie de página ---
\pagestyle{fancy}
\fancyhf{}
\fancyhead[L]{Basic C Without Tears}
\fancyhead[C]{Guía Especial -- El Arte del Pseudocódigo}
\fancyhead[R]{Stack}
\fancyfoot[L]{Basic-C-Without-Tears}
\fancyfoot[R]{\thepage}
\renewcommand{\headrulewidth}{0.4pt}
\renewcommand{\footrulewidth}{0.4pt}

% --- Entornos de Código C y Pseudocódigo ---
\lstdefinestyle{customc}{
  language=C,
  numbers=left,
  stepnumber=1,
  numbersep=6pt,
  numberstyle=\tiny\color{gray},
  basicstyle=\footnotesize\ttfamily,
  keywordstyle=\bfseries\color{blue!80!black},
  commentstyle=\itshape\color{green!50!black},
  stringstyle=\color{red!70!black},
  breaklines=true,
  tabsize=4,
  showstringspaces=false,
  frame=none
}

\lstdefinestyle{pseudocodigo}{
  numbers=left,
  stepnumber=1,
  numbersep=6pt,
  numberstyle=\tiny\color{gray},
  basicstyle=\footnotesize\ttfamily,
  keywordstyle=\bfseries\color{purple!80!black},
  commentstyle=\itshape\color{gray},
  morekeywords={INICIO, FIN, SI, ENTONCES, SINO, FIN_SI, MIENTRAS, HACER, FIN_MIENTRAS, PARA, DESDE, HASTA, FIN_PARA, REPETIR, HASTA_QUE, RETORNAR, LEER, ESCRIBIR, FUNCION, FIN_FUNCION, VERDADERO, FALSO},
  breaklines=true,
  tabsize=4,
  showstringspaces=false,
  frame=none
}

% --- Cajas estilizadas para comparación ---
\newtcolorbox{cbox}[1]{
  colback=blue!3,
  colframe=blue!70!black,
  arc=1mm,
  boxrule=0.6pt,
  title=\textbf{#1},
  coltitle=white,
  fonttitle=\small\sffamily,
  left=2mm, right=2mm, top=1.5mm, bottom=1.5mm
}

\newtcolorbox{pseudobox}[1]{
  colback=purple!3,
  colframe=purple!70!black,
  arc=1mm,
  boxrule=0.6pt,
  title=\textbf{#1},
  coltitle=white,
  fonttitle=\small\sffamily,
  left=2mm, right=2mm, top=1.5mm, bottom=1.5mm
}

\newtcolorbox{conceptbox}[1]{
  colback=gray!5,
  colframe=black,
  arc=0mm,
  boxrule=0.6pt,
  title=\textbf{#1},
  coltitle=white,
  fonttitle=\small\sffamily,
  left=2.5mm, right=2.5mm, top=1.5mm, bottom=1.5mm
}

\begin{document}

% =============================================================================
% PÁGINA 1: FUNDAMENTOS Y TABLA DE EQUIVALENCIAS
% =============================================================================
\begin{center}
    {\LARGE \textbf{Basic C Without Tears}} \\[0.2cm]
    {\Large \textbf{Guía de Referencia: El Arte del Pseudocódigo en C}} \\[0.1cm]
    \small Conceptos, Estándares de Escritura y Equivalencias con Código Compilable \\[0.05cm]
    \textbf{Autor:} Stack
\end{center}

\vspace{0.1cm}

\begin{conceptbox}{1. ¿Qué es el Pseudocódigo y cuándo utilizarlo?}
\small
El \textbf{pseudocódigo} es una descripción estructurada y rigurosa de un algoritmo. No depende de la sintaxis estricta de un compilador (puntos y coma, cabeceras o tipos de formateo), pero \textbf{no es lenguaje coloquial informal}.
\begin{itemize}[noitemsep, topsep=1pt]
    \item \textbf{Diseño previo:} Permite estructurar la solución antes de lidiar con la memoria y punteros en C.
    \item \textbf{Evaluaciones formativas y entrevistas:} Permite evaluar el razonamiento algorítmico y el control de flujo sin penalizar por detalles sintácticos menores de bibliotecas externas.
    \item \textbf{Regla de Oro:} Todo bloque debe tener delimitadores explícitos (\texttt{SI... FIN\_SI}, \texttt{MIENTRAS... FIN\_MIENTRAS}) y respetar sangría visual (\textit{indentación}).
\end{itemize}
\end{conceptbox}

\vspace{0.15cm}

\begin{conceptbox}{2. Tabla de Equivalencias Rápidas: Sintaxis C vs. Pseudocódigo Estándar}
\centering
\footnotesize
\begin{tabular}{p{4.2cm} p{5.5cm} p{5.5cm}}
\toprule
\textbf{Operación} & \textbf{Código Compilable en C} & \textbf{Pseudocódigo Válido} \\
\midrule
Asignación de variable & \texttt{total = 0;} & \texttt{total $\leftarrow$ 0} \quad ó \quad \texttt{total = 0} \\
Lectura de datos & \texttt{scanf("\%d", \&n);} & \texttt{LEER n} \\
Impresión en consola & \texttt{printf("Res: \%d\textbackslash n", x);} & \texttt{ESCRIBIR "Res: ", x} \\
Condicional simple & \texttt{if (cond) \{ ... \}} & \texttt{SI cond ENTONCES ... FIN\_SI} \\
Condicional doble & \texttt{if (c) \{ ... \} else \{ ... \}} & \texttt{SI c ENTONCES ... SINO ... FIN\_SI} \\
Bucle con contador & \texttt{for (int i=0; i<n; i++)} & \texttt{PARA i $\leftarrow$ 0 HASTA n-1 HACER} \\
Bucle condicional previo & \texttt{while (cond) \{ ... \}} & \texttt{MIENTRAS cond HACER ... FIN\_MIENTRAS} \\
Bucle con post-condición & \texttt{do \{ ... \} while (cond);} & \texttt{REPETIR ... HASTA QUE cond} \\
Acceso a arreglos & \texttt{A[i] = A[i - 1];} & \texttt{A[i] $\leftarrow$ A[i - 1]} \\
Retorno de valor & \texttt{return resultado;} & \texttt{RETORNAR resultado} \\
\bottomrule
\end{tabular}
\end{conceptbox}

% =============================================================================
% PÁGINA 2: COMPARACIÓN 1 - CONDICIONALES Y VALIDACIÓN
% =============================================================================
\newpage
\noindent {\Large \textbf{3. Comparativas Prácticas: Código C vs. Pseudocódigo}} \\[0.2cm]

\noindent \textbf{Caso 1: Validación y Clasificación de Datos (Control Condicional)} \\
\small \textit{Problema: Recibir tres medidas flotantes, verificar si conforman un triángulo admisible y clasificarlo por sus lados.}
\vspace{0.2cm}

\begin{cbox}{Código Compilable en C (C99/C11)}
\begin{lstlisting}[style=customc]
#include <stdio.h>

void clasificar_triangulo(float a, float b, float c) {
    // 1. Validacion de desigualdad triangular
    if (a < b + c && b < a + c && c < a + b) {
        printf("Triangulo valido: ");
        // 2. Clasificacion por igualdad de lados
        if (a == b && b == c) {
            printf("Equilatero\n");
        } else if (a == b || b == c || a == c) {
            printf("Isosceles\n");
        } else {
            printf("Escaleno\n");
        }
    } else {
        printf("No conforma un triangulo.\n");
    }
}
\end{lstlisting}
\end{cbox}

\vspace{0.25cm}

\begin{pseudobox}{Pseudocódigo Válido y Riguroso}
\begin{lstlisting}[style=pseudocodigo]
FUNCION clasificar_triangulo(a, b, c)
    // 1. Validar la desigualdad triangular
    SI (a < b + c) Y (b < a + c) Y (c < a + b) ENTONCES
        ESCRIBIR "Triangulo valido: "
        
        // 2. Clasificacion por igualdad de lados
        SI (a = b) Y (b = c) ENTONCES
            ESCRIBIR "Equilatero"
        SINO SI (a = b) O (b = c) O (a = c) ENTONCES
            ESCRIBIR "Isosceles"
        SINO
            ESCRIBIR "Escaleno"
        FIN_SI
    SINO
        ESCRIBIR "No conforma un triangulo."
    FIN_SI
FIN_FUNCION
\end{lstlisting}
\end{pseudobox}

\vspace{0.2cm}
\begin{conceptbox}{Clave Didáctica del Caso 1}
\small
En el pseudocódigo nos enfocamos en la lógica booleana (\texttt{Y}, \texttt{O}) y el anidamiento limpio. No es necesario gestionar bibliotecas (\texttt{\#include <stdio.h>}), prototipos ni secuencias de escape (\texttt{\textbackslash n}).
\end{conceptbox}

% =============================================================================
% PÁGINA 3: COMPARACIÓN 2 - ARREGLOS Y DETECCIÓN DE PICOS
% =============================================================================
\newpage
\noindent \textbf{Caso 2: Búsqueda y Filtrado en Arreglos Unidimensionales} \\
\small \textit{Problema: Recorrer un arreglo $A$ de $n$ elementos y almacenar en $picos$ aquellos valores estrictamente mayores a sus dos vecinos contiguos ($1 \le i \le n-2$).}
\vspace{0.2cm}

\begin{cbox}{Código Compilable en C (C99/C11)}
\begin{lstlisting}[style=customc]
#include <stdio.h>

int filtrar_picos(const int A[], int n, int picos[]) {
    int total_picos = 0;

    // Recorrido seguro de los indices interiores
    for (int i = 1; i < n - 1; i++) {
        if (A[i] > A[i - 1] && A[i] > A[i + 1]) {
            picos[total_picos] = A[i];
            total_picos++;
        }
    }
    return total_picos; // Retorna cantidad detectada
}
\end{lstlisting}
\end{cbox}

\vspace{0.25cm}

\begin{pseudobox}{Pseudocódigo Válido y Riguroso}
\begin{lstlisting}[style=pseudocodigo]
FUNCION filtrar_picos(A, n, picos)
    total_picos <- 0

    // Iterar estrictamente sobre elementos interiores
    PARA i <- 1 HASTA n - 2 HACER
        SI (A[i] > A[i - 1]) Y (A[i] > A[i + 1]) ENTONCES
            picos[total_picos] <- A[i]
            total_picos <- total_picos + 1
        FIN_SI
    FIN_PARA

    RETORNAR total_picos
FIN_FUNCION
\end{lstlisting}
\end{pseudobox}

\vspace{0.2cm}
\begin{conceptbox}{Clave Didáctica del Caso 2}
\small
Nota cómo el pseudocódigo respeta con precisión matemática los \textbf{límites del arreglo}: el ciclo recorre desde \texttt{1} hasta \texttt{n - 2}. Esta exactitud en los índices distingue a un pseudocódigo universitario sólido de una descripción informal vaga.
\end{conceptbox}

% =============================================================================
% PÁGINA 4: COMPARACIÓN 3 - STRINGS Y ERRORES COMUNES
% =============================================================================
\newpage
\noindent \textbf{Caso 3: Cadenas de Caracteres y Detección de Centinela (\texttt{'\textbackslash 0'})} \\
\small \textit{Problema: Determinar la longitud de una cadena recorriendo carácter por carácter hasta hallar el terminador nulo.}
\vspace{0.15cm}

\begin{cbox}{Código Compilable en C (C99/C11)}
\begin{lstlisting}[style=customc]
int longitud_cadena(const char str[]) {
    int len = 0;
    while (str[len] != '\0') {
        len++;
    }
    return len;
}
\end{lstlisting}
\end{cbox}

\vspace{0.15cm}

\begin{pseudobox}{Pseudocódigo Válido y Riguroso}
\begin{lstlisting}[style=pseudocodigo]
FUNCION longitud_cadena(str)
    len <- 0
    MIENTRAS str[len] != '\0' HACER
        len <- len + 1
    FIN_MIENTRAS
    RETORNAR len
FIN_FUNCION
\end{lstlisting}
\end{pseudobox}

\vspace{0.2cm}
\begin{conceptbox}{4. Criterios de Validez: Lo que SÍ se permite vs. Lo que INVALIDA una respuesta}
\small
\begin{tabular}{p{7.3cm} p{7.3cm}}
\toprule
\textbf{\textcolor{green!50!black}{[VÁLIDO] Prácticas Permitidas}} & \textbf{\textcolor{red!70!black}{[INVÁLIDO] Errores Graves}} \\
\midrule
Omitir directivas de compilación (\texttt{\#include}) y tipos de formato (\texttt{\%d}, \texttt{\%f}). & Escribir frases ambiguas: \textit{<<Buscar el mayor elemento>>} sin escribir el bucle de búsqueda. \\
\addlinespace[2pt]
Usar operadores intuitivos: $\leftarrow$ o \texttt{=} para asignar, y palabras en español como \texttt{Y}, \texttt{O}, \texttt{NO}. & Inventar funciones mágicas que resuelvan el núcleo del problema algorítmico pedido. \\
\addlinespace[2pt]
Omitir llaves \texttt{\{\}} siempre que la indentación y palabras de cierre delimiten el bloque. & Dejar bucles sin condición explícita de parada o variables acumuladoras sin inicializar. \\
\addlinespace[2pt]
Indexar arreglos y cadenas de texto mediante corchetes familiares \texttt{A[i]} o \texttt{str[k]}. & Omitir los cierres de estructuras (\texttt{FIN\_SI}, \texttt{FIN\_MIENTRAS}), volviendo ambiguo el anidamiento. \\
\bottomrule
\end{tabular}
\end{conceptbox}

\end{document}
